\section{AI × Space y nuevos mercados: de los sensores a las decisiones sobre recursos}

La discusión sobre space en la entrevista abarca communication, Earth observation, navigation, agricultura, agua y lunar robotics. Si solo se enumeran estos proyectos uno junto a otro, el resultado es un catálogo promocional; una estructura con mayor valor didáctico es la sensing-to-decision chain: la space infrastructure genera datos, la AI los interpreta, el gobierno o las empresas modifican la physical action y, después, los resultados se usan para calibrar el modelo.

\subsection{Communication, observation y harsh environment}

Esta sección reorganiza los proyectos de space vistos en clase como una sensing-to-decision pipeline. Al leer la figura, primero distingue las capacidades básicas que aportan la orbit y el sensor; después sigue cómo la imagery, el positioning y la connectivity llegan a la AI interpretation; por último, verifica si farming, water allocation o rover action reciben retroalimentación desde tierra. Solo cuando el último paso cierra el ciclo, los space data dejan de ser un producto informativo y se convierten en capacidad de gestión de recursos.

\lecturefigure{14-ai-space-stack.png}{La cadena de sistemas entre AI y space se extiende desde la infraestructura de percepción hasta las decisiones sobre recursos y las decisiones físicas.}{Subtítulos locales 24:00--27:10, 48:50--49:20; figura redibujada a partir del concepto.}

\begin{knowledgebox}{Lectura de la figura: la dificultad de la Earth observation está en el label y la intervention}
Las imágenes satelitales permiten identificar land use, vegetación o stress hídrico, pero la salida del modelo solo cambia los resultados si entra en un workflow de inspection, permit, irrigation o enforcement. Si la verdad de terreno es escasa, si la nubosidad y las variaciones estacionales son grandes y si las intervenciones de política pública no quedan registradas, se rompe la relación entre la precisión del modelo y el ahorro real de agua.
\end{knowledgebox}

\begin{warningbox}{Las proporciones de ahorro de agua y económicas mencionadas en clase son cifras del ponente}
Los subtítulos ofrecen varias proporciones sobre desertification, water saving, navigation y space-economy, pero carecen de un método y una línea base comunes. Este texto no toma esas cifras como conclusiones cuantitativas; solo conserva el patrón de sistema según el cual «geospatial sensing + AI + policy action» requiere una validación de ciclo cerrado.
\end{warningbox}

\subsection{NEOM y testbed: los nuevos mercados no eximen de la exigencia de evidencia}

Una ciudad nueva o un proyecto de gran escala puede reducir la legacy integration y ofrecer estándares unificados de identity, sensor e infrastructure, por lo que resulta adecuado como testbed; sin embargo, el carácter «greenfield» también amplifica los riesgos de scope, capital, governance y adoption. La aceptación del sistema debe partir igualmente de SLO de servicio en un ámbito reducido, unit economics, safety incidents y resident outcomes, en lugar de sustituir los datos de operación con la visión del proyecto.

\begin{importantbox}{Los cuatro niveles de evidencia de un testbed}
\begin{enumerate}
  \item el prototype funciona;
  \item el pilot es estable con una población y un entorno controlados;
  \item la production se sostiene con carga real, fallas y restricciones organizacionales;
  \item la transfer puede reproducirse en distintas ciudades, marcos regulatorios y condiciones de proveedores.
\end{enumerate}
Solo al alcanzar los dos últimos niveles queda demostrado que las capacidades de la plataforma son transferibles.
\end{importantbox}

\subsection{Resumen de la sección}

Lo esencial de AI × space no es la lista de satélites o rovers, sino si sensing, interpretation, decision y outcome forman un ciclo cerrado. Un greenfield testbed puede acelerar la experimentación, pero no reduce la exigencia de evidencia en producción ni de capacidad de transferencia.
